\documentclass[12pt]{report}
% Template Tesi in Informatica, ispirata al template "Tesi di laurea - Universit� di Pisa" di Simone Schirinzi (https://www.overleaf.com/latex/templates/tesi-di-laurea-universita-di-pisa/rwdcqtqwftpg). 

% Carattere dimensione 12

% Per la stampa fronte-retro sostituire con:
% \documentclass[12pt, twoside]{report}

% Margini (4cm a sx, 2.5cm a dx, 2.5cm in alto, 2.5cm in basso)
\usepackage[top=2.5cm, bottom=2.5cm, left=4cm, right=2.5cm, head=21.75pt, centering]{geometry}

% Per la stampa fronte-retro sostituire con: 
% \usepackage[top=2.5cm, bottom=2.5cm, inner=4cm, outer=4cm, right=2.5cm, centering]{geometry}

% Interlinea
\linespread{1.5}

% Librerie utili
\usepackage[italian]{babel} % applicazione regole di scrittura per la lingua italiana 
\usepackage[utf8]{inputenc} % codifica UTF-8
\usepackage{scrlayer-scrpage} % stili pagina per il frontespizio
\ifoot[]{}
\cfoot[]{}
\ofoot[\pagemark]{\pagemark}
\pagestyle{scrplain}
\usepackage{mathptmx} % font Times New Roman (simile)
\usepackage{graphicx} % inserimento di immagini
\usepackage{csquotes} % per le citazioni "in blocco"
%\usepackage[backend=biber, sorting=nty, ]{biblatex} % bibliografia con pacchetto biblatex (https://ctan.org/pkg/biblatex?lang=en)
\appto{\bibsetup}{\raggedright}

\usepackage{titlesec} % per la formattazione dei titoli delle sezioni, capitoli etc.
\usepackage{float} % per il posizionamento delle immagini

\usepackage{listings} % per il codice di programmazione
% Fonte https://en.wikibooks.org/wiki/LaTeX/Source_Code_Listings. Per la lista di sintassi riconosciute.
\renewcommand{\lstlistingname}{Code}% Listing -> Codice
\usepackage{xcolor}  % stile del codice
\definecolor{mygreen}{rgb}{0,0.6,0}
\definecolor{mygray}{rgb}{0.5,0.5,0.5}
\definecolor{mymauve}{rgb}{0.58,0,0.82}
\definecolor{darkgray}{rgb}{.4,.4,.4}
\definecolor{navy}{HTML}{000080}
\definecolor{purple}{rgb}{0.65, 0.12, 0.82}
\definecolor{codepurple}{rgb}{0.58,0,0.82}
\definecolor{backcolour}{rgb}{0.95,0.95,0.92}

% Stili configurabili del codice (lslisting) 
\lstset{ %
  belowcaptionskip=0.5em,
  backgroundcolor=\color{backcolour}, % choose the background color; you must add \usepackage{color} or \usepackage{xcolor}
  basicstyle=\footnotesize, % the size of the fonts that are used for the code
  breakatwhitespace=false, % sets if automatic breaks should only happen at whitespace
  breaklines=true, % sets automatic line breaking
  captionpos=b, % sets the caption-position to bottom
  commentstyle=\color{mygreen}, % comment style
  deletekeywords={...}, % if you want to delete keywords from the given language
  escapeinside={\%*}{*)}, % if you want to add LaTeX within your code
  extendedchars=true, % lets you use non-ASCII characters; for 8-bits encodings only, does not work with UTF-8
  frame=single, % adds a frame around the code
  keepspaces=true, % keeps spaces in text, useful for keeping indentation of code (possibly needs columns=flexible)
  keywordstyle=\color{codepurple}, % keyword style
  language=Python, % the language of the code
  morekeywords={*,...}, % if you want to add more keywords to the set
  numbers=left, % where to put the line-numbers; possible values are (none, left, right)
  numbersep=5pt, % how far the line-numbers are from the code
  numberstyle=\tiny\color{mygray}, % the style that is used for the line-numbers
  rulecolor=\color{black}, % if not set, the frame-color may be changed on line-breaks within not-black text (e.g. comments (green here))
  showspaces=false, % show spaces everywhere adding particular underscores; it overrides 'showstringspaces'
  showstringspaces=false, % underline spaces within strings only
  showtabs=false, % show tabs within strings adding particular underscores
  stepnumber=1, % the step between two line-numbers. If it's 1, each line will be numbered
  stringstyle=\color{mymauve}, % string literal style
  tabsize=2, % sets default tabsize to 2 spaces
  title=\lstname % show the filename of files included with \lstinputlisting; also try caption instead of title
}

% END of listing package 

% Esempio riconoscimento sintassi JavaScript 

\lstdefinelanguage{JavaScript}{
  keywords={typeof, new, true, false, catch, function, return, null, catch, switch, var, const, let, if, in, while, do, else, case, break},
  keywordstyle=\color{blue}\bfseries,
  ndkeywords={class, export, boolean, throw, implements, import, this},
  ndkeywordstyle=\color{darkgray}\bfseries,
  identifierstyle=\color{black},
  sensitive=false,
  comment=[l]{//},
  morecomment=[s]{/*}{*/},
  commentstyle=\color{mygreen}\ttfamily,
  stringstyle=\color{red}\ttfamily,
  morestring=[b]',
  morestring=[b]"
}

\lstset{
   language=JavaScript,
   extendedchars=true,
   basicstyle=\footnotesize\ttfamily,
   showstringspaces=false,
   showspaces=false,
   numbers=left,
   numberstyle=\footnotesize,
   numbersep=9pt,
   tabsize=2,
   breaklines=true,
   showtabs=false,
   captionpos=b
}

% Formato delle intestazioni
\titleformat{\chapter}[block]
  {\normalfont\LARGE\bfseries}{\thechapter.}{0.5em}{\LARGE}
\titlespacing*{\chapter}{0pt}{-20pt}{25pt}

\begin{document}

% Frontespizio
\begin{titlepage}
  \begin{figure}
    \centering\includegraphics[scale=0.5]{immagini/cherubino_pant541.png}
  \end{figure}

  \begin{center}
    {\LARGE{ Corso di Laurea in Informatica \\ }}
    \vspace{2cm}
    {\Large { TESI DI LAUREA }}\\
    \vspace{2cm}
    {\Large { Algoritmi genetici per Explainable AI }}
  \end{center}

  \vspace{2cm}

  \begin{minipage}[t]{0.47\textwidth}
    {\large{\bf Relatori:\\ Riccardo Guidotti}}
    \vspace{0.5cm}
    {\large{\bf \\Correlatore:\\ Marco Danelutto}}
  \end{minipage}\hfill\begin{minipage}[t]{0.47\textwidth}\raggedleft
    {\large{\bf Candidato: \\ Federico Bustaffa\\ }}
  \end{minipage}

  \vspace{25mm}

  \centering{\large{\bf ANNO ACCADEMICO 2024/2025 }}
\end{titlepage}
% Fine frontespizio

\tableofcontents
\thispagestyle{empty}

\listoffigures

\thispagestyle{empty}
\clearpage
\setcounter{page}{1}
\addtocontents{toc}{\protect\thispagestyle{empty}}
\addcontentsline{toc}{chapter}{Introduzione} % Capitolo non numerato
\chapter*{Introduzione}
\label{ch:introduzione}

Introduzione...

\clearpage

\chapter{Capitolo 1}
\label{ch:capitolo1}

% --- Inizio del Capitolo 1 ---

Capitolo 1...

Esempio di una nota\footnote{CleanCode} con citazione a sezione~\ref{sec:sezione1}.

\begin{figure}[H]
    \centering
    \includegraphics[width=0.2\textwidth]{immagini/logo-dip_blu_hr.png}
    \caption{Esempio di un'immagine}
    \label{fig:immagine1}
\end{figure}

\begin{lstlisting}[caption={Esempio codice JavaScript},label={Esempio codice JavaScript}, language=JavaScript]
function foo(num) {
    const bar = 2;
    return num + bar;
}

const result = foo(2);
// result -> 4
\end{lstlisting}


\section{Sezione 1\label{sec:sezione1}}


Sezione 1... con citazione bibliografica~\cite{CleanCode}

\subsection{Sottosezione 1}
\label{subsec:sottosezione1}

Sottosezione 1...

\clearpage

\addcontentsline{toc}{chapter}{Conclusioni} % Capitolo non numerato
\chapter*{Conclusioni}
\label{ch:conclusioni}

Conclusioni...

\bibliographystyle{plain} % We choose the "plain" reference style
\bibliography{bibliography} % Entries are in the refs.bib file
%\addbibresource{bibliography.bib}
%\printbibliography

\end{document}